% Ad Atticum 11.5 --- headnote
% ad-atticum-11-05, 4 November 48 BC, Brundisium

\textit{Cicero to Atticus, written from Brundisium on the day
before the Nones of November 48 BC --- 4 November (the
manuscript dateline: \textit{Scr.\ Brundisi pr.\ Non.\ Nov.\
a.\ 706 (48)}; \textit{works.yaml} currently records 5 November
for this letter, but \textit{pridie Nonas} is 4 November). The
political ground has shifted entirely since 11.4: Pompey is
dead, killed on the Egyptian shore on 28 September, and Cicero
has come back to Italy on the strength of his right to use the
\textit{imperium} of a returning proconsul. He has landed at
Brundisium without the cause and without his colleagues, and
sits there now exposed --- to Caesar's pleasure, to the resentment
of the remaining Pompeians, and to the suspicion of Antony, who
holds Italy in Caesar's absence.}

\textit{The letter is in four short sections. \S1 is the apology
he will give again and again over the next year for having
quit the Pompeian camp: the causes were so bitter and so heavy
and so unprecedented that he acted ``by a kind of impulse of
the mind rather than by deliberation,'' and they produced what
Atticus can now see. He understands from Atticus's letters
that Atticus is hunting for new ways to protect him. \S2
declines, gently, Atticus's suggestion that he come closer to
Rome travelling by night through the towns: he has no lodgings
suitable for hiding by day, and for the purpose he is wanted
it hardly matters whether he is seen on the road or in a town.
\S3 explains why he has written so little: distress of mind
and body has kept him from finishing letters; he asks Atticus
to write in his name to Basilus, Servilius, and anyone else
who seems fit. \S4 answers a query about Vatinius (any
goodwill is welcome, if anyone can find a way to use it), then
reports that his brother Quintus met him at Patrae in a thoroughly
hostile frame of mind, that his nephew joined Quintus there
from Corcyra, and that the two of them must have set out from
Patrae with the rest of the Pompeian remnant. The rupture with
Quintus will dominate the rest of Book 11.}
